\section{Tensore energia impulso}
Come prima applicazione del teorema di Noether consideriamo l'invarianza della Lagrangiana per traslazioni nello spazio tempo per il campo di Klein-Gordon.
$$
\mathcal{L} = \frac{1}{2}[\partial^\mu\phi\partial_\mu\phi - m^2\phi^2]
$$
Le traslazioni sono elementi del gruppo di Poincarré e si tratta di trasformazioni dello spazio-tempo. Una traslazione è definita come
$$
T_a : x \to x + a \qquad\text{ per le componenti }\qquad  x^\nu \to x^\nu + a^\nu
$$
Trasformazioni differenziabili che dipendono da $4$ parametri $a^\nu$. Per una traslazione infinitesima $\delta a$ ($\delta a$ piccoli) i campi hanno una variazione
$$
\delta\phi(x) = \phi(x+\delta a) - \phi(x) \approx \phi(x) + \partial_\nu\phi\delta a^\nu - \phi(x) = \delta a^\nu\partial_\nu\phi \qquad \frac{\delta\phi}{\delta a^\nu}=\partial_\nu\phi
$$
La corrente di Noether è pertanto
$$
j_{\mu\nu} = \frac{\partial\mathcal{L}}{\partial(\partial^\mu\phi)}\frac{\delta\phi}{\delta a^\nu} \qquad j_{\mu\nu} = \partial_\mu\phi\partial_\nu\phi
$$
Abbiamo visto che
$$
\delta\mathcal{L} = \partial_\mu \sum_n \frac{\partial\mathcal{L}}{\partial(\partial_\mu\phi_n)}\delta\phi_n
$$
Ma attenzione $\delta \mathcal{L}$ non è nulla. In definitiva
$$
\delta\mathcal{L} = \partial^\mu j_{\mu\nu}\delta a^\nu\qquad\frac{\delta\mathcal{L}}{\delta a^\nu} = \partial^\mu j_{\mu\nu} = \partial^\mu \partial_\mu\phi\partial_\nu\phi
$$
Infatti la lagrangiana è implicitamente una funzione della posizione. La variazione indotta dalla traslazione è
$$
\mathcal{L}(x) = \mathcal{L}(\phi(x), \partial_\mu\phi(x)) \qquad \delta\mathcal{L} = \delta a^\mu \partial_\mu \mathcal{L}=g_{\mu\nu}\partial^\mu\delta a^\nu\mathcal{L}(x)
$$
$$
\frac{\partial\mathcal{L}}{\partial a^\nu} = \partial^\mu g_{\mu\nu} \mathcal{L}(x)
$$
Uguagliando con l'analoga espressione trovata per la variazione dei campi 
$$
\partial^\mu \partial_\mu\phi\partial_\nu\phi = \partial^\mu g_{\mu\nu}\mathcal{L}(x)
$$
Definiamo il tensore Energia-Impulso per il campo di Klein-Gordon
$$
T_{\mu\nu} = \partial_\mu\phi\partial_\nu\phi - g_{\mu\nu}\mathcal{L}(x)
$$
La conservazione della corrente è $\partial^\mu T_{\mu\nu}=0$. Le "cariche" conservate sono le componenti del $4-$momento.
$$
P_\nu=\int T_{0\nu}d^3\mathbf{r}
$$
In particolare l'Hamiltoniana
$$
H \equiv P_0 = \int T_{00} d^3\mathbf{r} \qquad \mathcal{H} = T_{00} = \partial_0\phi\partial_0\phi - g_{00}\mathcal{L}(x) \qquad \mathcal{H} = \dot{\phi}\dot{\phi} - \mathcal{L}(x)
$$
Analizziamo anche le componenti $\nu=1,3$ del $4-$momento
$$
T_{\mu\nu} = \partial_\mu\phi\partial_\nu\phi - g_{\mu\nu}\mathcal{L}(x)
$$
$$
P^\nu = \int T_0^\nu d^3\mathbf{r} = \int \partial_0\hat{\phi}\partial^\nu\hat{\phi} d^3\mathbf{r} = -\int \hat{\dot{\phi}}\nabla\hat{\phi} d^3\mathbf{r}
$$
Anche nel caso di $P$, come per $H$, alla fine avremo bisogno di eliminare un contributo infinito. Oppure adottiamo l'ordinamento normale.
$$
\hat{\mathbf{P}} = -\int : \hat{\dot{\phi}}\nabla\hat{\phi} : d^3\mathbf{r}
$$
Sviluppiamo l'espressione
$$
\mathbf{P} = -\int : \hat{\dot{\phi}}\boldsymbol{\nabla}\hat{\phi} : d^3\mathbf{r} = -\frac{1}{2} \int : \hat{\dot{\phi}}\boldsymbol{\nabla}\hat{\phi} + \hat{\dot{\phi}}\boldsymbol{\nabla}\hat{\phi} : d^3\mathbf{r}
$$
Osserviamo che 
$$
\hat{\dot{\phi}}\boldsymbol{\nabla}\hat{\phi} = \boldsymbol{\nabla}(\hat{\dot{\phi}}\hat{\phi}) - (\boldsymbol{\nabla}\hat{\dot{\phi}})\hat{\phi} \qquad \int \partial_i(\hat{\dot{\phi}}\hat{\phi}) d^3\mathbf{r} = \int dx_j dx_k \int_{-\infty}^{+\infty} \partial_i(\hat{\dot{\phi}}\hat{\phi}) dx_i \to 0
$$
Inoltre $:(\boldsymbol{\nabla}\hat{\dot{\phi}})\hat{\phi}: \quad = \quad :\hat{\phi}(\boldsymbol{\nabla}\hat{\dot{\phi}}):$. Ovviamente le due espressione conducono a risultati uguale solo se si applicano gli operatori al vuoto $|0\rangle$ quindi otteniamo
$$
\mathbf{P} = -\frac{1}{2} \int : \hat{\dot{\phi}}\boldsymbol{\nabla}\hat{\phi} - \hat{\phi}\boldsymbol{\nabla}\hat{\dot{\phi}} : d^3\mathbf{r}
$$
A questo punto si sostituiscono le rppresentazioni del campo e si trova, come per l'Hamiltoniano
$$
\mathbf{P} = \frac{1}{(2\pi)^3} \int \mathbf{p} \hat{a}_\mathbf{p}^\dagger \hat{a}_\mathbf{p} d^3\mathbf{p}
$$
\section{Il campo scalare complesso}
Il campo di Klein-Gordon che abbiamo fino ad ora studiato ha solo un grado di libertà. Infatti si utilizza per descrivere particelle neutre senza spin: bosoni neutri, particelle che coincidono con la propria antiparticella. Per descrivere particelle cariche abbiamo bisogno di un campo complesso e la carica va intesa in senso ampio. Infatti è un numero quantico che distingue particelle e antiparticelle. Non è necessariamente la carica elettrica. Consideriamo quindi due particelle di Klein-Gordon con la stessa massa. La Lagrangiana di questo sistema è
$$
\mathcal{L}(x) = \frac{1}{2}\partial_\mu\hat{\phi}_1\partial^\mu\hat{\phi}_1 - \frac{1}{2}m^2\hat{\phi}_1^2 + \frac{1}{2}\partial_\mu\hat{\phi}_2\partial^\mu\hat{\phi}_2 - \frac{1}{2}m^2\hat{\phi}_2^2
$$
I campi $\phi_1$ e $\phi_2$ sono due gradi di libertà di un sistema che ha bisogno di due componenti in uno spazio astratto. Studiamo gli effetti dell'invarianza rispetto a rotazioni in questo spazio.
$$
\hat{\phi}'_1 = \hat{\phi}_1 \cos\alpha - \hat{\phi}_2 \sin\alpha
$$
$$
\hat{\phi}'_2 = \hat{\phi}_1 \sin\alpha + \hat{\phi}_2 \cos\alpha
$$
Le trasformazioni dipendono da un parametro (gruppo $SO(2)$). Consideriamo una rotazione infinitesima: $\alpha\to\delta\alpha$ quindi $\cos{\delta\alpha}\approx0$ e $\sin{\delta\alpha}\approx\delta\alpha$. Quindi la trasformazione e quindi la variazione dei due campi sono
$$
\hat{\phi}'_1 = \hat{\phi}_1 - \hat{\phi}_2\delta\alpha \qquad \delta\hat{\phi}_1 = \hat{\phi}'_1 - \hat{\phi}_1 = -\hat{\phi}_2\delta\alpha \qquad \frac{\delta\hat{\phi}_1}{\delta\alpha} = -\hat{\phi}_2
$$
$$
\hat{\phi}'_2 = \hat{\phi}_1\delta\alpha + \hat{\phi}_2 \qquad \delta\hat{\phi}_2 = \hat{\phi}'_2 - \hat{\phi}_2 = \hat{\phi}_1\delta\alpha \qquad \frac{\delta\hat{\phi}_2}{\delta\alpha} = \hat{\phi}_1
$$
Si tratta di una simmetria interna con un solo parametro. La corrente di Noether è
\begin{align*}
\hat{N}^\mu_\phi &= \sum_n \frac{\partial\mathcal{L}}{\partial(\partial_\mu\hat{\phi}_n)} \frac{\delta\hat{\phi}_n}{\delta\alpha} = \frac{\partial\mathcal{L}}{\partial(\partial_\mu\hat{\phi}_1)} \frac{\delta\hat{\phi}_1}{\delta\alpha} + \frac{\partial\mathcal{L}}{\partial(\partial_\mu\hat{\phi}_2)} \frac{\delta\hat{\phi}_2}{\delta\alpha} \\
&= (\partial^\mu\hat{\phi}_1)(-\hat{\phi}_2) + (\partial^\mu\hat{\phi}_2)\hat{\phi}_1
\end{align*}
$$
\hat{N}^\mu_\phi = (\partial^\mu\hat{\phi}_2)\hat{\phi}_1 - (\partial^\mu\hat{\phi}_1)\hat{\phi}_2
$$
La carica conservata (da confrontare con la corrente dell'equazione di Klein-Gordon) è
$$
\hat{Q} = \int \hat{N}^0_\phi d^3\mathbf{r} = \int \left( \dot{\hat{\phi}}_2\hat{\phi}_1 - \dot{\hat{\phi}}_1\hat{\phi}_2 \right) d^3\mathbf{r}
$$
Risulta conveniente introdurre un campo complesso $\phi$
$$
\hat{\phi} = \frac{1}{\sqrt{2}}(\hat{\phi}_1 - i\hat{\phi}_2) \qquad \hat{\phi}^\dagger = \frac{1}{\sqrt{2}}(\hat{\phi}_1 + i\hat{\phi}_2)
$$
Ricordando lo sviluppo di un campo scalare reale
$$
\hat{\phi}_n(x) = \frac{1}{(2\pi)^3} \int \frac{d^3\mathbf{k}}{\sqrt{2\omega}} \left[ \hat{a}_{n\mathbf{k}} e^{-ik\cdot x} + \hat{a}_{n\mathbf{k}}^\dagger e^{+ik\cdot x} \right]
$$
Troviamo il seguente sviluppo per il campo complesso
$$
\hat{\phi}(x) = \frac{1}{(2\pi)^3} \int \frac{d^3\mathbf{k}}{\sqrt{2\omega}} \left[ \hat{a}_{\mathbf{k}} e^{-ik\cdot x} + \hat{b}_{\mathbf{k}}^\dagger e^{+ik\cdot x} \right]
$$
Gli operatore $a$ e $b^\dagger$ sono adesso indipendenti
$$
\hat{a}_{\mathbf{k}} = \frac{1}{\sqrt{2}}(\hat{a}_{1\mathbf{k}} - i\hat{a}_{2\mathbf{k}}) \qquad \hat{b}_{\mathbf{k}}^\dagger = \frac{1}{\sqrt{2}}(\hat{a}_{1\mathbf{k}}^\dagger - i\hat{a}_{2\mathbf{k}}^\dagger)
$$
Si può verificare che gli operatori di creazione e distruzione introdotti hanno le seguenti regole di commutazione
$$
[\hat{a}_{\mathbf{k}}, \hat{a}_{\mathbf{k}'}^\dagger] = (2\pi)^3 \delta(\mathbf{k} - \mathbf{k}') \qquad [\hat{b}_{\mathbf{k}}, \hat{b}_{\mathbf{k}'}^\dagger] = (2\pi)^3 \delta(\mathbf{k} - \mathbf{k}')
$$
Tutte le altre regole di commutazione sono nulle. Si può verificare facilmente che nel piano complesso la trasformazione che abbiamo introdotto precedentemente diventa
$$
\hat{\phi}' = e^{-i\alpha}\hat{\phi} \qquad \hat{\phi}' = (1 - i\delta\alpha)\hat{\phi} \qquad \delta\hat{\phi} = -i\delta\alpha\hat{\phi} \qquad \delta\hat{\phi}^\dagger = i\delta\alpha\hat{\phi}^\dagger
$$
La trasformazione in oggetto prende il nome di trasformazione di fase (gruppo $U(1)$) o trasformazione di gauge globale (la fase non dipende da $x$). La Lagrangiana per il campo complesso è 
$$
\mathcal{L} = \partial_\mu\hat{\phi}^\dagger\partial^\mu\hat{\phi} - m^2\hat{\phi}^\dagger\hat{\phi}
$$
Ripetiamo il calcolo precedente utilizzando il campo complesso e quindi la corrente di Noether è
\begin{align*}
\hat{N}^\mu_\phi &= \sum_n \frac{\partial\mathcal{L}}{\partial(\partial_\mu\hat{\phi}_n)} \frac{\delta\hat{\phi}_n}{\delta\alpha} = \frac{\partial\mathcal{L}}{\partial(\partial_\mu\hat{\phi}^\dagger)} \frac{\delta\hat{\phi}^\dagger}{\delta\alpha} + \frac{\partial\mathcal{L}}{\partial(\partial_\mu\hat{\phi})} \frac{\delta\hat{\phi}}{\delta\alpha} \\
&= (\partial^\mu\hat{\phi})(i\hat{\phi}^\dagger) + (\partial^\mu\hat{\phi}^\dagger)(-i\hat{\phi}) \qquad \hat{N}^\mu_\phi = i \left[ (\partial^\mu\hat{\phi})\hat{\phi}^\dagger - (\partial^\mu\hat{\phi}^\dagger)\hat{\phi} \right]
\end{align*}
La carica conservata è (da confrontare con la corrente dell'equazione di Klein-Gordon)
$$
\hat{Q} = \int \hat{N}^0_\phi d^3\mathbf{r} = i \int \left[ \dot{\hat{\phi}}\hat{\phi}^\dagger - \dot{\hat{\phi}}^\dagger\hat{\phi} \right] d^3\mathbf{r}
$$
Utilizzando le espansioni integrali si possono calcolare l'Hamiltoniana
$$
\hat{H} = \frac{1}{(2\pi)^3} \int d^3\mathbf{k} [\hat{a}_{\mathbf{k}}^\dagger \hat{a}_{\mathbf{k}} + \hat{b}_{\mathbf{k}}^\dagger \hat{b}_{\mathbf{k}}] \omega
$$
e l'operatore Carica
$$
\hat{Q} = \frac{1}{(2\pi)^3} \int d^3\mathbf{k} [\hat{a}_{\mathbf{k}}^\dagger \hat{a}_{\mathbf{k}} - \hat{b}_{\mathbf{k}}^\dagger \hat{b}_{\mathbf{k}}]
$$
In entrambi i calcoli si sono trascurati i termini che danno contributo infinito (equivalente a imporre prodotti normali). Entrambi gli operatori ($H$ e $Q$) si possono esprimere in funzione di due operatori numero relativi a due differenti tipi di quanti
$$
\hat{n}_{\mathbf{k}}^+ = \hat{a}_{\mathbf{k}}^\dagger \hat{a}_{\mathbf{k}} \qquad \hat{n}_{\mathbf{k}}^- = \hat{b}_{\mathbf{k}}^\dagger \hat{b}_{\mathbf{k}}
$$
Entrambi gli operatori numero ($n^+$ e $n^-$) hanno autovalori non negativi. Visto che l'Hamiltoniana contiene la somma di questi operatori allora è definita positiva e si è eliminato il problema delle energie negative. L'operatore Carica è la differenza dei due operatori. I due quanti contribuiscono con segno opposto alla carica. Hanno carica opposta (sono uno l'antiparticella dell'altro). Analizziamo un po' meglio l'effetto degli operatori $a$ e $b$. Per un generico operatore di creazione o distruzione $o_p$
$$
[\hat{Q}, \hat{O}_\mathbf{p}] = \frac{1}{(2\pi)^3} \int d^3\mathbf{k} [\hat{a}_\mathbf{k}^\dagger \hat{a}_\mathbf{k} - \hat{b}_\mathbf{k}^\dagger \hat{b}_\mathbf{k}, \hat{O}_\mathbf{p}]
$$
Per i vari operatori di creazione e distruzione si ha
$$
[\hat{a}_\mathbf{k}^\dagger \hat{a}_\mathbf{k} - \hat{b}_\mathbf{k}^\dagger \hat{b}_\mathbf{k}, \hat{a}_\mathbf{p}] = [\hat{a}_\mathbf{k}^\dagger \hat{a}_\mathbf{k}, \hat{a}_\mathbf{p}] = \hat{a}_\mathbf{k}^\dagger [\hat{a}_\mathbf{k}, \hat{a}_\mathbf{p}] + [\hat{a}_\mathbf{k}^\dagger, \hat{a}_\mathbf{p}] \hat{a}_\mathbf{k} = - [\hat{a}_\mathbf{p}, \hat{a}_\mathbf{k}^\dagger] \hat{a}_\mathbf{k} = -(2\pi)^3 \delta(\mathbf{k} - \mathbf{p}) \hat{a}_\mathbf{k}
$$
Pertanto
$$
\begin{aligned}
[\hat{a}_\mathbf{k}^\dagger \hat{a}_\mathbf{k} - \hat{b}_\mathbf{k}^\dagger \hat{b}_\mathbf{k}, \hat{a}_\mathbf{p}] &= -(2\pi)^3 \delta(\mathbf{k} - \mathbf{p}) \hat{a}_\mathbf{k} \\
[\hat{a}_\mathbf{k}^\dagger \hat{a}_\mathbf{k} - \hat{b}_\mathbf{k}^\dagger \hat{b}_\mathbf{k}, \hat{a}_\mathbf{p}^\dagger] &= (2\pi)^3 \delta(\mathbf{k} - \mathbf{p}) \hat{a}_\mathbf{p}^\dagger \\
[\hat{a}_\mathbf{k}^\dagger \hat{a}_\mathbf{k} - \hat{b}_\mathbf{k}^\dagger \hat{b}_\mathbf{k}, \hat{b}_\mathbf{p}] &= (2\pi)^3 \delta(\mathbf{k} - \mathbf{p}) \hat{b}_\mathbf{p} \\
[\hat{a}_\mathbf{k}^\dagger \hat{a}_\mathbf{k} - \hat{b}_\mathbf{k}^\dagger \hat{b}_\mathbf{k}, \hat{b}_\mathbf{p}^\dagger] &= -(2\pi)^3 \delta(\mathbf{k} - \mathbf{p}) \hat{b}_\mathbf{p}^\dagger
\end{aligned}
$$
Inserendo nella relazione di partenza
$$
[\hat{Q}, \hat{a}_\mathbf{p}] = -\hat{a}_\mathbf{p} \qquad [\hat{Q}, \hat{a}_\mathbf{p}^\dagger] = \hat{a}_\mathbf{p}^\dagger \qquad [\hat{Q}, \hat{b}_\mathbf{p}] = \hat{b}_\mathbf{p} \qquad [\hat{Q}, \hat{b}_\mathbf{p}^\dagger] = -\hat{b}_\mathbf{p}^\dagger
$$
\begin{itemize}
    \item Gli operatori $a$ e $b^\dagger$ diminuiscono di $1$ la carica: l'operatore $a$ distrugge una particella, l'operatore $b^\dagger$ crea un'antiparticella
    \item Gli operatori $a^\dagger$ e $b$ aumentano di $1$ la carica: l'operatore $a^\dagger$ crea una particella, l'operatore $b$ distrugge un'antiparticella.
\end{itemize}
Gli stati del campo complesso si costruiscono come nel caso del campo reale. Adesso però ci sono due tipi di particelle indipendenti
\begin{itemize}
    \item Particelle "Positive" $|\mathbf{p},+\rangle=\hat{a}_\mathbf{p}^\dagger|0\rangle$
    \item Particelle "Negative" $|\mathbf{p},-\rangle=\hat{b}_\mathbf{p}^\dagger|0\rangle$
\end{itemize}
Analogamente si costruiscono stati con un numero arbitrario di quanti dei due tipi. Data la forma degli operatori $Q$ e $H$ è immediato verificare che l'operatore $Q$ è conservato:
$$
\dot{\hat{Q}} = [\hat{Q}, \hat{H}] = 0
$$
In assenza di interazioni si può facilmente verificare che il numero di antiparticelle e il numero di particelle si conservano separatamente. Si può inoltre verificare che per i campi $\phi$ e $\phi^\dagger$ valgono le regole di commutazione con $Q$ simili a quelle degli operatori di creazione e distruzione
$$
[\hat{Q}, \hat{\phi}] = -\hat{\phi} \qquad [\hat{Q}, \hat{\phi}^\dagger] = \hat{\phi}^\dagger
$$
Definiamo tramite sviluppo in serie l'operatore 
$$
\hat{U}(\alpha) = e^{i\alpha\hat{Q}}\quad \text{ si può dimostrare che }\quad \hat{U}(\alpha)\hat{\phi}\hat{U}^{-1}(\alpha) = e^{i\alpha}\hat{\phi}
$$
\section{Quantizzazione del Campo di Dirac}
La quantizzazione del campo di Dirac $\psi$ si effettua in modo del tutto analogo alla quantizzazione del campo di Klein-Gordon. Seguendo il nostro approccio per il campo $KG$:
\begin{itemize}
    \item si sviluppa il campo nei modi normali
    \item si promuovono i coefficienti dello sviluppo a operatori
    \item si impongono le regole di commutazioni
\end{itemize}
Alternativamente si può seguire il metodo canonico:
\begin{itemize}
    \item si scrive la Lagrangiana del campo di Dirac
    \item si individuano i momenti coniugati ai campi
    \item si impongono le relazioni di commutazioni fra i campi e momenti coniugati
\end{itemize}
Anticicipiamo la differenza fondamentale sta nel fatto che nelle regole di commutazione ai commutatori $[A,B]$ sostituiamo gli anticommutatori $\{A,B\}$. La necessità degli anticommutatori può essere motivata alternativamente dal principio di esclusione di Pauli che richiede l'antisimmetria degli stati fermionici costruiti con gli operatori di creazione e distruzione. Inoltre anche dalla necessità che l'Hamiltoniana sia definita positiva. Entrambe le impostazioni derivano da motivazioni fisiche e vale la pena approfondire entrambi gli aspetti.
\section{Lagrangiana del campo di Dirac}
La Lagrangiana del campo di Dirac è
$$
\mathcal{L}(\psi, \psi^\dagger, \partial_\mu\psi) = \overline{\psi}(i\slashed{\partial} - m)\psi
$$
Notiamo che non contiene i campo $\partial_\mu^\dagger$. L'equazione di Dirac segue semplicemente applicando le equazioni di Eulero Lagrange ai campi $\psi$ e $\psi^\dagger$. 
$$
\partial_\mu \frac{\partial\mathcal{L}}{\partial(\partial_\mu\psi_n)} = \frac{\partial\mathcal{L}}{\partial\psi_n}
$$
Variando i campi $\phi^\dagger$
$$
\partial_\mu \frac{\partial\mathcal{L}}{\partial(\partial_\mu\psi^\dagger)} = \frac{\partial\mathcal{L}}{\partial\psi^\dagger} \qquad 0 = \gamma^0(i\slashed{\partial} - m)\psi \qquad (i\slashed{\partial} - m)\psi = 0
$$
Alternativamente, variando i campi $\psi$
$$
\partial_\mu \frac{\partial\mathcal{L}}{\partial(\partial_\mu\psi)} = \frac{\partial\mathcal{L}}{\partial\psi} \qquad \partial_\mu i\overline{\psi}\gamma^\mu = -m\overline{\psi} \qquad \overline{\psi}(i\overleftarrow{\slashed{\partial}} + m) = 0
$$
L'asimmetria tra $\psi$ e $\psi^\dagger$ deriva dal fatto che più Lagrangiane sono possibili, ad esempio
$$
\mathcal{L}_1 = \mathcal{L}^* = \overline{[\overline{\psi}(i\slashed{\partial} - m)\psi]} = \overline{[\overline{\psi}(i\gamma^\mu\partial_\mu\psi - m\psi)]}=\left(\overline{i\gamma^\mu\partial_\mu\psi}-m\overline{\psi}\right)\psi
$$
$$
\mathcal{L}_1 = -\overline{\psi}(i\overleftarrow{\slashed{\partial}} + m)\psi
$$
Tuttavia si può dimostrare che le azioni corrispondenti differiscono per l'integrale di una $4-$divergenza che si può rendere nullo all'infinito. Pertanto conducono alle stesse equazioni. Pertanto conducono alle stesse equazioni.
\section{Hamiltoniana del campo di Dirac}
Possiamo a questo punto calcolare i momenti coniugati $\psi$. Ovviamente i momenti coniugati $\psi^\dagger$ sono nulli per l'asimmetria tra $\psi$ e $\psi^\dagger$ nella Lagrangiana: $\mathcal{L} = \overline{\psi}(i\slashed{\partial} - m)\psi$
$$
\pi = \frac{\partial\mathcal{L}}{\partial(\partial_0\psi)} = \overline{\psi}i\gamma^0 = i\psi^\dagger\gamma^0\gamma^0 = i\psi^\dagger
$$
L'Hamiltoniana è pertanto
\begin{align*}
    \mathcal{H} &= \pi\partial_0\psi - \mathcal{L} = i\psi^\dagger\partial_0\psi - \overline{\psi}(i\slashed{\partial} - m)\psi = i\psi^\dagger\gamma^0\gamma^0\partial_0\psi - \psi^\dagger\gamma^0(i\gamma^\mu\partial_\mu - m)\psi \\
    &= -\psi^\dagger\gamma^0 i \left( \sum_{k=1,3} \gamma^k\partial_k - m \right) \psi = \psi^\dagger (-i\boldsymbol{\alpha} \cdot \boldsymbol{\nabla} + \gamma^0 m) \psi \qquad \mathcal{H} = \psi^\dagger (-i\boldsymbol{\alpha} \cdot \boldsymbol{\nabla} + m\beta) \psi
\end{align*}
Dal momento che $\psi$ soddisfa l'equazione di Dirac: $(-i\boldsymbol{\alpha} \cdot \boldsymbol{\nabla} + m\beta)\psi = i\partial_0\psi$ abbiamo una utile espressione alternativa per l'Hamiltoniana $H=\psi^\dagger i \partial_0\psi$. Si potrebbe calcolare il momento del campo calcolando il tensore $T^{\mu\nu}$. Diamo i due risultati (il calcolo può essere un esercizio)
$$
H = \int d^3\mathbf{r} \psi^\dagger i\partial_0\psi \qquad \mathbf{P} = \int d^3\mathbf{r} \psi^\dagger (-i\boldsymbol{\nabla})\psi
$$
\section{Espansione del campo di Dirac}
L'espansione del campo di Dirac in onde piane si può fare allo stesso modo di quanto fatto per il campo scalare
$$
\hat{\psi}(x) = \int \sum_{\sigma=\pm s} \frac{d^3\mathbf{k}}{(2\pi)^3\sqrt{2E_k}} \left( \hat{a}_{k,\sigma} u_{k,\sigma} e^{-ik\cdot x} + \hat{b}_{k,\sigma}^\dagger v_{k,\sigma} e^{+ik\cdot x} \right)
$$
Conviene definire delle funzioni con ben definite proprietà di ortogonalità
$$
f_{k,\sigma} = \frac{u_{k,\sigma} e^{-ik\cdot x}}{(2\pi)^3\sqrt{2E_k}} \qquad g_{k,\sigma} = \frac{v_{k,\sigma} e^{+ik\cdot x}}{(2\pi)^3\sqrt{2E_k}}
$$
Si possono verificare
$$
\int f_{k\sigma}^\dagger(x) f_{k'\sigma'}(x) d^3\mathbf{r} = \delta_{\sigma,\sigma'}\delta_{k,k'} \qquad \int g_{k\sigma}^\dagger(x) g_{k'\sigma'}(x) d^3\mathbf{r} = \delta_{\sigma,\sigma'}\delta_{k,k'}
$$
$$
\int f_{k\sigma}^\dagger(x) g_{k'\sigma'}(x) d^3\mathbf{r} = 0 \qquad \int g_{k\sigma}^\dagger(x) f_{k'\sigma'}(x) d^3\mathbf{r} = 0
$$
Il campo diventa 
$$
\hat{\psi}(x) = \int \sum_{\sigma=\pm s} d^3\mathbf{k} (f_{k,\sigma} \hat{a}_{k,\sigma} + g_{k,\sigma} \hat{b}_{k,\sigma}^\dagger)
$$
Utilizzando queste relazioni si può facilmente inveritre l'espansione
$$
\hat{a}_{p\sigma} = \int f_{p\sigma}^\dagger(x) \hat{\psi}(x) d^3\mathbf{r} \qquad \hat{b}_{p\sigma}^\dagger = \int g_{p\sigma}^\dagger(x) \hat{\psi}(x) d^3\mathbf{r}
$$
\section{Quantizzazione del campo di Dirac II}
Possiamo adesso discutere la quantizzazione del campo di Dirac. Seguendo il nostro approccio iniziale promuoviamo le funzioni $a_{k\sigma}$ e $b_{k\sigma}$ a operatori in uno spazio astratto e fissiamo le regole di anti-commutazione
$$
\{\hat{a}_{k\sigma}, \hat{a}_{k'\sigma'}^\dagger\} = (2\pi)^3 \delta_{\sigma\sigma'}\delta_{kk'} \qquad \{\hat{b}_{k\sigma}, \hat{b}_{k'\sigma'}^\dagger\} = (2\pi)^3 \delta_{\sigma\sigma'}\delta_{kk'} \qquad \text{tutte le altre sono nulle}
$$
Con questi operatori si possono costruire gli stati con $n$ particelle nello spazio di Fock
$$
|k_1\sigma_1, \dots, k_n\sigma_n\rangle = \sqrt{2E_{k_1}}\dots\sqrt{2E_{k_n}} \, \hat{a}_{k_1\sigma_1}^\dagger \dots \hat{a}_{k_n\sigma_n}^\dagger |0\rangle
$$
Una analoga espressione si può scrivere con gli operatori $b^\dagger$ (antiparticelle). Le regole di anti-commutazione tra gli operatori $a$ e $b$ assicurano che gli stati hanno la corretta antisimmetria rispetto alle permutazioni di due particelle.
$$
|k_1\sigma_1, \dots, \textcolor{red}{k_j\sigma_j}, \dots,\textcolor{red}{k_i\sigma_i}, \dots, k_n\sigma_n\rangle = \pm |k_1\sigma_1, \dots, \textcolor{red}{k_i\sigma_i}, \dots, \textcolor{red}{k_j\sigma_j}, \dots, k_n\sigma_n\rangle
$$
Il segno è $+$ o $-$ per permutazioni pari o dispari rispettivamente. Nel seguito sarà utile trattare separatamente le due componenti di un campo
$$
\hat{\psi}(x) = \hat{\psi}_a(x) + \hat{\psi}_b(x)
$$
\begin{itemize}
    \item La componente a energia positiva
    $$
\hat{\psi}_a(x) = \int \sum_{\sigma=\pm s} f_{k\sigma}(x) \hat{a}_{k,\sigma} d^3k
$$
\item La componente a energia negativa
$$
\hat{\psi}_b^\dagger(x) = \int \sum_{\sigma=\pm s} g_{k\sigma}(x) \hat{b}_{k,\sigma}^\dagger d^3k
$$
\end{itemize}
Studiamo l'effetto del campo $\psi(x)$ sul vuoto
$$
\hat{\psi}(x) = \int \sum_{\sigma=\pm s} \frac{d^3\mathbf{k}}{(2\pi)^3\sqrt{2E_k}} \left( \hat{a}_{k,\sigma} u_{k,\sigma} e^{-ik\cdot x} + \hat{b}_{k,\sigma}^\dagger v_{k,\sigma} e^{+ik\cdot x} \right)
$$
In particolare l'effetto di $\psi^\dagger(x)$
\begin{align*}
\hat{\psi}^\dagger |0\rangle &= \int \sum_{\sigma=\pm s} \frac{d^3\mathbf{k}}{(2\pi)^3\sqrt{2E_k}} \left( \hat{a}_{k,\sigma}^\dagger u_{k,\sigma}^\dagger e^{+ik\cdot x} |0\rangle + \hat{b}_{k,\sigma} v_{k,\sigma}^\dagger e^{-ik\cdot x} |0\rangle \right) \\
&= \int \sum_{\sigma=\pm s} \frac{d^3\mathbf{k}}{(2\pi)^3\sqrt{2E_k}} \hat{a}_{k,\sigma}^\dagger u_{k,\sigma}^\dagger e^{+ik\cdot x} |0\rangle
\end{align*}
Notiamo che $\hat{a}_{k,\sigma}^\dagger|0\rangle$ è uno stato di una particella di momento $\mathbf{k}$. Pertanto $\hat{\psi}|0\rangle$ è la sovrapposizione di infiniti stati con infiniti valori di $\mathbf{k}$. Effettuiamo l'analisi dello stato come nel caso di Klein-Gordon. I due stati sono
$$
|\mathbf{p}, \rho, a\rangle = \sqrt{2E_p} \hat{a}_{p\rho}^\dagger |0\rangle \qquad (\hat{\psi}^\dagger |0\rangle)^\dagger = \langle 0 | \hat{\psi} = \int \sum_{\sigma=\pm s} \frac{d^3\mathbf{k}}{(2\pi)^3\sqrt{2E_k}} \langle 0 | \hat{a}_{k,\sigma} u_{k,\sigma} e^{-ik\cdot x}
$$
Otteniamo
\begin{align*}
\langle 0 | \hat{\psi} | \mathbf{p}, \rho, a \rangle &= \int \sum_{\sigma=\pm s} \frac{d^3\mathbf{k}}{(2\pi)^3} \frac{\sqrt{2E_p}}{\sqrt{2E_k}} \langle 0 | \hat{a}_{k,\sigma} u_{k,\sigma} e^{-ik\cdot x} \hat{a}_{p\rho}^\dagger | 0 \rangle \\
&= \int \sum_{\sigma=\pm s} \frac{d^3\mathbf{k}}{(2\pi)^3} \frac{\sqrt{2E_p}u_{k,\sigma} e^{-ik\cdot x}}{\sqrt{2E_k}} \langle 0 | \hat{a}_{k,\sigma} \hat{a}_{p\rho}^\dagger | 0 \rangle
\end{align*}
Ricordiamo che
$$
\{\hat{a}_{k\sigma}, \hat{a}_{k'\sigma'}^\dagger\} = (2\pi)^3 \delta_{\sigma\sigma'}\delta_{kk'}
$$
Otteniamo
\begin{align*}
\langle 0 | \hat{\psi} | \mathbf{p}, \rho, a \rangle &= \int \sum_{\sigma=\pm s} \frac{d^3\mathbf{k}}{(2\pi)^3} \frac{\sqrt{2E_p}}{\sqrt{2E_k}} u_{k,\sigma} e^{-ik\cdot x} \langle 0 | \hat{a}_{p,\rho}^\dagger \hat{a}_{k,\sigma} + (2\pi)^3 \delta_{kp} \delta_{\rho\sigma} | 0 \rangle \\
&= \int \sum_{\sigma=\pm s} \frac{d^3\mathbf{k}}{(2\pi)^3} \frac{u_{k,\sigma} e^{-ik\cdot x} \sqrt{2E_p}}{\sqrt{2E_k}} (2\pi)^3 \delta_{kp} \delta_{\rho\sigma} = u_{p\rho} e^{-ik\cdot x}
\end{align*}
Ancora una volta abbimo trovato una soluzione dell'equazione di Dirc: onda piana di una particella localizzata in $x$. Analogamente si dimostr che $\psi(x)$ produce uno stato di antiparticella. Con un calcolo simile al precedente si può dimostrare che
$$
\langle \mathbf{p}, \rho, b | \hat{\psi} | 0 \rangle = v_{p\rho} e^{+ip\cdot x}
$$
Con altri conti si può dimostrare che $\psi(x)$ produce uno stato di antiparticella svolgendo analoghi passaggi ma usando l'anticommutazione di $b$ e $b^\dagger$ (saltabile). Inoltre è possibile anche verificare le regole di anticommutazione tra gli $\hat{\psi}$ a tempi uguali (non saltabile ma conti complessi) ottenendo 
$$
\{\hat{\psi}_n(x), \hat{\psi}_m^\dagger(x')\} = \delta_{nm} \delta(\mathbf{r} - \mathbf{r}')
$$
$$
\{\hat{\psi}_n(x), \hat{\psi}_m(x')\} = 0 \qquad \{\hat{\psi}_n^\dagger(x), \hat{\psi}_m^\dagger(x')\} = 0
$$
\section{Hamiltoniana e regole di commutazioni}
Veniamo adesso all'espressione dell'Hamiltoniana in funzione degli operatori di creazione e distruzione. Per questo è conveniente partire dalle relazioni
$$
H = \int d^3\mathbf{r} \hat{\psi}^\dagger i \partial_0 \hat{\psi} \qquad \hat{\psi}(x) = \int \sum_{\sigma=\pm s} d^3\mathbf{k} (f_{k,\sigma} \hat{a}_{k,\sigma} + g_{k,\sigma} \hat{b}_{k,\sigma}^\dagger)
$$
Utilizzando l'espansione di $\psi$ otteniamo
$$
i\partial_0 \hat{\psi}(x) = \int \sum_{\sigma'=\pm s} d^3\mathbf{k}' E_{k'} (f_{k',\sigma'} \hat{a}_{k',\sigma'} - g_{k',\sigma'} \hat{b}_{k',\sigma'}^\dagger) \qquad \hat{\psi}^\dagger(x) = \int \sum_{\sigma=\pm s} d^3\mathbf{k} (f_{k,\sigma}^\dagger \hat{a}_{k,\sigma}^\dagger + g_{k,\sigma}^\dagger \hat{b}_{k,\sigma})
$$
Introducendo nell'Hamiltoniana
$$
H = \int d^3\mathbf{r} \hat{\psi}^\dagger i \partial_0 \hat{\psi} = \int d^3\mathbf{r} \sum_{\sigma,\sigma'=\pm s} d^3\mathbf{k} d^3\mathbf{k}' (f_{k,\sigma}^\dagger \hat{a}_{k,\sigma}^\dagger + g_{k,\sigma}^\dagger \hat{b}_{k,\sigma}) E_{k'} (f_{k',\sigma'} \hat{a}_{k',\sigma'} - g_{k',\sigma'} \hat{b}_{k',\sigma'}^\dagger)
$$
Eseguiamo l'integrazione sul volume e utilizziamo le relazioni di ortogonalità: sopravvivono solo due termini
$$
H = \int \sum_{\sigma,\sigma'=\pm s} d^3\mathbf{k} d^3\mathbf{k}' E_{k'} \int d^3\mathbf{r} (f_{k,\sigma}^\dagger f_{k',\sigma'} \hat{a}_{k,\sigma}^\dagger \hat{a}_{k',\sigma'} - g_{k,\sigma}^\dagger g_{k',\sigma'} \hat{b}_{k,\sigma} \hat{b}_{k',\sigma'}^\dagger)
$$
$$
= \int \sum_{\sigma,\sigma'=\pm s} d^3\mathbf{k} d^3\mathbf{k}' E_{k'} (\hat{a}_{k,\sigma}^\dagger \hat{a}_{k',\sigma'} \delta_{\sigma\sigma'}\delta_{kk'} - \hat{b}_{k,\sigma} \hat{b}_{k',\sigma'}^\dagger \delta_{\sigma\sigma'}\delta_{kk'}) = \int \sum_{\sigma=\pm s} d^3\mathbf{k} E_k (\hat{a}_{k,\sigma}^\dagger \hat{a}_{k,\sigma} - \hat{b}_{k,\sigma} \hat{b}_{k,\sigma}^\dagger)
$$
Il termine relativo ai quanti $a$ è già sottoforma di un operatore numero. Per mettere anche il secondo termine sotto forma di operatore numero abbiamo bisogno di commutare i due operatori. Se usassimo una regola di commutazione (trascuriamo il termine infinito)
$$
[\hat{b}_{k,\sigma}, \hat{b}_{k,\sigma}^\dagger] = c \to \hat{b}_{k,\sigma} \hat{b}_{k,\sigma}^\dagger = c + \hat{b}_{k,\sigma}^\dagger \hat{b}_{k,\sigma} \qquad H = \int \sum_{\sigma=\pm s} d^3k E_k (\hat{a}_{k,\sigma}^\dagger \hat{a}_{k,\sigma} - \hat{b}_{k,\sigma} \hat{b}_{k,\sigma}^\dagger)
$$
L'Hamiltoniana non sarebbe definita positiva. I quanti $b$ contribuirebbero con energia negativa. Al contrario, con una regola di anti-commutazione
$$
\{\hat{b}_{k,\sigma}, \hat{b}_{k,\sigma}^\dagger\} = c \to \hat{b}_{k,\sigma} \hat{b}_{k,\sigma}^\dagger = c - \hat{b}_{k,\sigma}^\dagger \hat{b}_{k,\sigma} \qquad \hat{H} = \int \sum_{\sigma=\pm s} d^3k E_k (\hat{a}_{k,\sigma}^\dagger \hat{a}_{k,\sigma} + \hat{b}_{k,\sigma}^\dagger \hat{b}_{k,\sigma})
$$
Questa volta anche i quanti $b$ contribuiscono all'energia con segno positivo. Adottare regole di anti-commutazione è necessario per avere una Hamiltoniana definita positiva.
\section{Invarianza gauge di globale}
Anche la Lagrangiana del campo di Dirac è invariante per trasformazioni di fase costante: trasformazioni di gauge globali
$$
\mathcal{L} = \overline{\psi}(i\slashed{\partial} - m)\psi \qquad \psi' = e^{-i\alpha}\psi \qquad \psi'^\dagger = e^{+i\alpha}\psi^\dagger
$$
Ovviamente, se la fase è costante, la Lagrangiana è invariante
$$
\mathcal{L}' = \overline{\psi}'(i\slashed{\partial} - m)\psi' = e^{+i\alpha}\overline{\psi}(i\slashed{\partial} - m)e^{-i\alpha}\psi = e^{+i\alpha}e^{-i\alpha}\overline{\psi}(i\slashed{\partial} - m)\psi = \mathcal{L}
$$
Calcoliamo adesso la corrente che corrisponde a questa invarianza, per una trasformazione infinitesima
$$
\psi' = (1 - i\delta\alpha)\psi \qquad \delta\psi = -i\delta\alpha\psi \qquad \frac{\delta\psi}{\delta\alpha} = -i\psi \quad \text{ analogamente }\quad\frac{\delta\psi^\dagger}{\delta\alpha} = i\psi^\dagger
$$
Ricordiamo l'espressione per la corrente di Noether
$$
j^\mu = \sum_n \frac{\partial\mathcal{L}}{\partial(\partial_\mu\psi_n)} \frac{\delta\psi_n}{\delta\alpha} = \frac{\partial\mathcal{L}}{\partial(\partial_\mu\psi)} \frac{\delta\psi}{\delta\alpha} + \cancel{\frac{\partial\mathcal{L}}{\partial(\partial_\mu\psi^\dagger)} \frac{\delta\psi^\dagger}{\delta\alpha}} \qquad \frac{\partial\mathcal{L}}{\partial(\partial_\mu\psi)} = i\overline{\psi}\gamma^\mu \qquad 
$$
$$
j^\mu = i\overline{\psi}\gamma^\mu (-i\psi)\qquad \Longrightarrow\qquad j^\mu = \overline{\psi}\gamma^\mu\psi \qquad \frac{\partial\mathcal{L}}{\partial(\partial_\mu\psi^\dagger)} = 0
$$
Abbiamo trovato la corrente di probabilità della teoria di Dirac che è una corrente conservata
$$
\partial_\mu j^\mu = 0
$$
Calcoliamo adesso la carica conservata associata alla conservazione della corrente
$$
\hat{Q} = \int j^0 d^3\mathbf{r} = \int \overline{\psi}(x)\gamma^0\psi(x) d^3\mathbf{r} = \int \psi^\dagger(x)\psi(x) d^3\mathbf{r}
$$
Le espansioni dei campi sono
$$
\hat{\psi}(x) = \int \sum_{\sigma=\pm s} d^3\mathbf{k} (f_{k,\sigma} \hat{a}_{k,\sigma} + g_{k,\sigma} \hat{b}_{k,\sigma}^\dagger) \qquad \hat{\psi}^\dagger(x) = \int \sum_{\sigma'=\pm s} d^3\mathbf{k}' (f_{k',\sigma'}^\dagger \hat{a}_{k',\sigma'}^\dagger + g_{k',\sigma'}^\dagger \hat{b}_{k',\sigma'})
$$
Inserendo le due espressioni nella carica
$$
\hat{Q} = \int d^3\mathbf{r} dk' dk \sum_{\sigma\sigma'} (f_{k',\sigma'}^\dagger \hat{a}_{k',\sigma'}^\dagger + g_{k',\sigma'}^\dagger \hat{b}_{k',\sigma'}) (f_{k,\sigma} \hat{a}_{k,\sigma} + g_{k,\sigma} \hat{b}_{k,\sigma}^\dagger)
$$
Utilizzando le proprietà di ortogonalità delle funzioni $f$ e $g$ sopravvivono solo due termini
$$
\hat{Q} = \int dk' dk \sum_{\sigma\sigma'} \int d^3\mathbf{r} (f_{k',\sigma'}^\dagger f_{k,\sigma} \hat{a}_{k',\sigma'}^\dagger \hat{a}_{k,\sigma} + g_{k',\sigma'}^\dagger g_{k,\sigma} \hat{b}_{k',\sigma'} \hat{b}_{k,\sigma}^\dagger)
$$
$$
\hat{Q} = \int dk' dk \sum_{\sigma\sigma'} (\hat{a}_{k',\sigma'}^\dagger \hat{a}_{k,\sigma} \delta_{\sigma\sigma'}\delta_{kk'} + \hat{b}_{k',\sigma'}^\dagger \hat{b}_{k,\sigma} \delta_{\sigma\sigma'}\delta_{kk'}) = \int dk \sum_{\sigma} (\hat{a}_{k,\sigma}^\dagger \hat{a}_{k,\sigma} + \hat{b}_{k,\sigma}^\dagger \hat{b}_{k,\sigma})
$$
vediamo che anche in questo caso per trasformare il secondo operatore in un operatore numero abbiamo bisogno delle regole di commutazione. Vediamo l'effetto di avere utilizzato anticommutatori
$$
\{\hat{b}_{k,\sigma}, \hat{b}_{k,\sigma}^\dagger\} = c \to \hat{b}_{k,\sigma} \hat{b}_{k,\sigma}^\dagger = c - \hat{b}_{k,\sigma}^\dagger \hat{b}_{k,\sigma} \qquad \hat{Q} = \int \sum_{\sigma=\pm s} d^3\mathbf{k} (\hat{a}_{k,\sigma}^\dagger \hat{a}_{k,\sigma} - \hat{b}_{k,\sigma}^\dagger \hat{b}_{k,\sigma})
$$
Con questa scelta le particelle e le antiparticelle hanno cariche opposte. Nei calcoli che abbiamo fatto ($Q$ e $H$) abbiamo scartato il contributo infinito che deriva dal commutatore. Nel caso dei campi fermionici abbiamo però cambiato anche il segno degli operatori. Dobbiamo precisare la definizione del prodotto normale nel caso di campi fermionici. IL prodotto normale porta gli operatori a distribuzioni a destra e inserisce un segno $+$ o $-$ a seconda del numero pari o dispari di permutazioni
$$
: \hat{a}_1 \hat{a}_2^\dagger : = - \hat{a}_2^\dagger \hat{a}_1 \qquad : \hat{b}_1 \hat{a}_1 \hat{a}_2^\dagger : = \hat{a}_2^\dagger \hat{b}_1 \hat{a}_1
$$
L'invarianza globale di gauge della Lagrangiana di Dirac porta alla corrente conservata
$$
\mathcal{L} = \overline{\psi}(i\slashed{\partial} - m)\psi \qquad j^\mu = \overline{\psi}\gamma^\mu\psi
$$
L'aspetto importante, conseguenza dell'invarianza di gauge, è il fatto che i campi appaiono nella forma
$$
\hat{j} = \overline{\psi}\hat{O}\psi
$$
\begin{itemize}
    \item Banalizzando, la presenza contemporanea di un campo hermitiano coniugato e un campo normale assicura che compaia il prodotto della fase e della sua complessa coniugata il cui prodotto è $1$
    \item Meno banale: la presenza contemporanea di un campo hermitiano coniugato e un campo "normale" assicura che la carica sia conservata. Pensiamo allo scattering. Particella che arriva particella che esce questa diventa una corrente e poi la leghiamo anche ad un fotone virtuale e andiamo a calcolare il valore di aspettazione di questo operatore di corrente. Questo operatore lo calcolo tra due stati ma se questi hanno una carica diversa questo valore è zero. Devono avere la stessa carica. Il fatto che ci sia un operatore aggiunto hermitiano e uno normale mi assicura che l'accoppiamento tra distruzione e costruzione mi conserva la carica.
\end{itemize}
I campi contengono operatori di creazione e distruzione: protoni e antiprotoni
$$
\hat{\psi}^\dagger(x) = \int \sum_{\sigma'=\pm s} d^3\mathbf{k} (f_{\mathbf{k},\sigma'}^\dagger \hat{a}_{\mathbf{k},\sigma'}^\dagger + g_{\mathbf{k},\sigma'}^\dagger \hat{b}_{\mathbf{k},\sigma'})
$$
$$
\hat{\psi}(x) = \int \sum_{\sigma=\pm s} d^3\mathbf{k} (f_{\mathbf{k},\sigma} \hat{a}_{\mathbf{k},\sigma} + g_{\mathbf{k},\sigma} \hat{b}_{\mathbf{k},\sigma}^\dagger)
$$
L'applicazione della corrente a uno stato induce a una variazione di carica nulla
$$
\Delta Q = 0
$$
\begin{center}
\begin{tabular}{|c|c|c|c|c|c|c|c|c|}
\hline
 & & \multicolumn{3}{c|}{$p$} & \multicolumn{3}{c|}{$\bar{p}$} & \\ 
\cline{1-9} 
$\psi^\dagger$ & $\psi$ & $\psi^\dagger$ & $\psi$ & $Q_p$ & $\psi^\dagger$ & $\psi$ & $Q_{\bar{p}}$ & $Q_T$\\ 
\hline
$a^\dagger$ & $a$ & $+1$ & $-1$ & $0$ & & & $0$ & $0$ \\ 
\hline
$a^\dagger$ & $b^\dagger$ & $+1$ & & $+1$ & & $+1$ & $-1$ & $0$ \\ 
\hline
$b$ & $a$ & & $-1$ & $-1$ & $-1$ & & $+1$ & $0$ \\ 
\hline
$b$ & $b^\dagger$ & & & $0$ & $-1$ & $+1$ & $0$ & $0$ \\ 
\hline
\end{tabular}
\end{center}
\section{Invarianza locale di gauge e QED}
Gli elettroni liberi sono descritti dalla Lagrangiana
$$
\mathcal{L}_e = i\overline{\psi}\gamma^\mu\partial_\mu\psi - m\overline{\psi}\psi
$$
Questa Lagrangiana è invariante per trasformazioni di gauge globali
\begin{itemize}
    \item Gruppo $U(1)$
    $$
\psi(x) \to \psi'(x) = e^{i\alpha}\psi(x)
$$
\end{itemize}
L'invarianza si perde se $\alpha$ dipende da $x$ (trasformazioni di gauge locali)
$$
\partial_\mu\psi'(x) = \partial_\mu e^{i\alpha(x)}\psi(x) = e^{i\alpha(x)}\partial_\mu\psi(x) \textcolor{red}{+ e^{i\alpha(x)}[i\partial_\mu\alpha(x)]\psi(x)}
$$
Il termine aggiuntivo, dovuto al fatto che $\alpha$ dipende da $x$, distrugge l'invarianza. La Lagrangiana modificata è
$$
\mathcal{L}'_e = \overline{\psi}(x)e^{-i\alpha(x)} e^{i\alpha(x)} \gamma^\mu \partial_\mu\psi(x) \textcolor{red}{+ \overline{\psi}(x)e^{-i\alpha(x)}e^{i\alpha(x)} \gamma^\mu [i\partial_\mu\alpha(x)]\psi(x)} - m\overline{\psi}(x)e^{-i\alpha(x)}e^{i\alpha(x)}\psi(x)
$$
$$
\mathcal{L}'_e = \mathcal{L}_e +\textcolor{red}{ \overline{\psi}(x) [\gamma^\mu i\partial_\mu\alpha(x)]\psi(x)} \qquad \text{non è invariante}
$$
Si può costruire una teoria in cui la Lagrangiana sia invariante per trasformazioni di gauge locali aggiungendo altri campi e modificando opportunamente l'operatore differenziale. Partiamo dal secondo punto, la modifica dell'operatore differenziale. Occorre compensare la derivata per la variazione del gauge (derivata covariante). Aggiungiamo un termine "moltiplicazione per funzione" all'operatore $\partial_\mu$.
$$
\partial_\mu \to D_\mu = \partial_\mu - ieA_\mu(x)
$$
$A_\mu$ serve per compensare il termine aggiuntivo introdotto dalla trasformazione locale di fase. Perchè questo avvenga è necessario che anche la funzione $A_\mu$ si trasformi in un modo ben determinato se si opera una trasformazione di gauge locale. Ricordiamo che il termine aggiuntivo
$$
\mathcal{L}'_e = \mathcal{L}_e + \overline{\psi}(x) [i\gamma^\mu \partial_\mu\alpha(x)]\psi(x)
$$
Per compensarlo
$$
A_\mu \to A'_\mu = A_\mu + \frac{1}{e}\partial_\mu\alpha(x)
$$
Ricordiamo l'elettrodinamica classica! Nella prossima diapositiva si verifica che la Lagrangiana scritta utilizzando la derivata covariante è invariante per le trasformazioni di gauge locali di $\psi$ e $A_\mu$. Infatti la nuova Lagrangiana è
$$
\mathcal{L} = i\overline{\psi}\gamma^\mu \textcolor{red}{D_\mu}\psi - m\overline{\psi}\psi = i\overline{\psi}\gamma^\mu (\partial_\mu - ieA_\mu)\psi - m\overline{\psi}\psi
$$
$$
\mathcal{L} = i\overline{\psi}\gamma^\mu\partial_\mu\psi - m\overline{\psi}\psi \textcolor{red}{+ e\overline{\psi}\gamma^\mu\psi A_\mu}
$$
Verifichaimoche la Lagrangiana scritta utilizzando la derivata covariante è invariante per le trasformazioni di gauge locali di $\psi$ e $A_\mu$. Abbiamo infatti visto nelle diapositive precedenti che
$$
i\overline{\psi}\gamma^\mu\partial_\mu\psi \to i\overline{\psi}\gamma^\mu\partial_\mu\psi \textcolor{red}{+ i\overline{\psi}\gamma^\mu [i\partial_\mu\alpha(x)]\psi} = i\overline{\psi}\gamma^\mu\partial_\mu\psi \textcolor{red}{- \overline{\psi}\gamma^\mu\psi\partial_\mu\alpha(x)}
$$
$$
e\overline{\psi}\gamma^\mu\psi A_\mu \to e\overline{\psi}\gamma^\mu\psi (A_\mu \textcolor{red}{+ \frac{1}{e}\partial_\mu\alpha(x))} = e\overline{\psi}\gamma^\mu\psi A_\mu\textcolor{red}{ + \overline{\psi}\gamma^\mu\psi\frac{1}{e}\partial_\mu\alpha(x)}
$$
I due termini aggiuntivi dovuto alle trasformazioni di gauge di $\psi$ e $A_\mu$ di elidono. Notiamo che il termine aggiuntivo introdotto dalla derivata covariante $D_\mu$ garantisce l'invarianza di gauge e rappresenta l'interazione dell'elettrone con un campo vettoriale. La teoria viene completata introducendo il termine cinetico per il campo $A_\mu$
$$
\mathcal{L}_\gamma = -\frac{1}{4} F^{\mu\nu} F_{\mu\nu} \qquad F_{\mu\nu} = \partial_\mu A_\nu - \partial_\nu A_\mu
$$
Il termine cinetico è la lagrangiana libera di una particella vetotriale di spin $1$ che è il campo elettromagnetico (il fotone). Questo termine per il campo vettoriale $A_\mu$ corrisponde al termine cinetico dell'elettrone $\overline{\psi}\gamma^\mu\partial_\mu\psi$ e descrive l'evoluzione liberla del campo $A_\mu$. Questa è la Lagrangiana del campo elettromagnetico. Si verifica facilmente che anche il nuovo termine cinetico possiede l'invarianza di gauge. In particolare $F^{\mu\nu}$ è già invariante. Se il campo $A_\mu$ descrive una particella con massa $m_\gamma$ la Lagrangiana dovrebbe contenere anche il termine
$$
\mathcal{L}_{m_\gamma} = m_\gamma A^\mu A_\mu
$$
Questo termine non sarebbe invariante rispetto alle trasformazioni di gauge del campo $A_\mu$. Per avere una teroia per campi vettoriale che possegga l'invarianza di gauge locale occorre che il campo $A_\mu$ sia una particella senza massa